% generator_I07.tex -- Prop I.7 (Theor.): unique triangle on base -- equal conterminous sides.
% Geometry and text from jemmybutton (CC-BY-SA). Build: lualatex generator_I07.tex
\documentclass[12pt]{article}
\usepackage{geometry}\geometry{a5paper, margin=1.5cm}
\usepackage[lining]{ebgaramond}
\usepackage{amsmath}
\usepackage{amssymb}
\usepackage{byrne}
\def\mpPre{u := 1cm; textLabels := false;}

\begin{document}

\defineNewPicture{%
  pair A, B, C, D, E, F, G, H;%
  A := (0, 0);%
  B := A shifted (4u, 0);%
  C := A shifted (u, 3u);%
  D := C shifted (7/4u, 0);%
  E := 1/2[C, D] yscaled -0.7;%
  F := E shifted (0, -2u);%
  G := 5/4[A, E];%
  H := 5/4[A, F];%
  byAngleDefine.C(B, C, A, byblack, SOLID_SECTOR);%
  byAngleDefine(D, C, B,  byred,    SOLID_SECTOR);%
  byAngleDefine.D(A, D, B, byyellow, SOLID_SECTOR);%
  byAngleDefine(C, D, A,  byblue,   SOLID_SECTOR);%
  byAngleDefine(B, F, H,  byblack,  SOLID_SECTOR);%
  byAngleDefine(B, F, E,  byred,    SOLID_SECTOR);%
  byAngleDefine(B, E, G,  byyellow, SOLID_SECTOR);%
  byAngleDefine(G, E, F,  byblue,   SOLID_SECTOR);%
  draw byNamedAngleResized();%
  draw byLine(C, D, byblack, DASHED_LINE, REGULAR_WIDTH);%
  draw byLine(E, F, byblack, DASHED_LINE, REGULAR_WIDTH);%
  draw byLine(A, B, byblack, SOLID_LINE,  REGULAR_WIDTH);%
  byLineDefine(B, C, byblue,  SOLID_LINE, REGULAR_WIDTH);%
  byLineDefine(C, A, byred,   SOLID_LINE, REGULAR_WIDTH);%
  byLineDefine(B, D, byblue,  SOLID_LINE, REGULAR_WIDTH);%
  byLineDefine(D, A, byred,   SOLID_LINE, REGULAR_WIDTH);%
  byLineDefine(B, E, byblue,  SOLID_LINE, REGULAR_WIDTH);%
  byLineDefine(E, A, byred,   SOLID_LINE, REGULAR_WIDTH);%
  byLineDefine(B, F, byblue,  SOLID_LINE, REGULAR_WIDTH);%
  byLineDefine(F, A, byred,   SOLID_LINE, REGULAR_WIDTH);%
  byLineDefine(E, G, byred,   DASHED_LINE, REGULAR_WIDTH);%
  byLineDefine(F, H, byred,   DASHED_LINE, REGULAR_WIDTH);%
  draw byNamedLine(EG,FH);%
  draw byNamedLineSeq(0)(BC,CA,EA,BE);%
  draw byNamedLineSeq(0)(BD,DA,FA,BF);%
  byPointLabelDefine(F, "C");%
  byPointLabelDefine(E, "D");%
  draw byLabelsOnPolygon(F, A, C, D, B, F, noPoint)(OMIT_FIRST_LABEL+OMIT_LAST_LABEL, 0);%
  draw byLabelsOnPolygon(A, E, B)(OMIT_FIRST_LABEL+OMIT_LAST_LABEL, 0);%
  draw byLabelsOnPolygon(H, F, A)(OMIT_FIRST_LABEL+OMIT_LAST_LABEL, 0);%
}

% STATEMENT
\noindent
\begin{minipage}[t]{0.45\linewidth}
\centering\vspace{0pt}%
\drawCurrentPicture
\end{minipage}\hfill
\begin{minipage}[t]{0.52\linewidth}
\vspace{0pt}%
\textbf{Book~I.\enspace Prop.\ VII. Theor.}

\medskip\noindent\itshape
On the same base (\drawUnitLine{AB}), and on the same side
of it there cannot be two triangles having their conterminous
sides (\drawUnitLine{CA} and \drawUnitLine{DA},
\drawUnitLine{BC} and \drawUnitLine{BD}) at both extremities
of the base, equal to each other.

\bigskip

% PROOF
\noindent\upshape
When two triangles stand on the same base, and on the same
side of it, the vertex of the one shall either fall outside
of the other triangle, or within it; or, lastly, on one of
its sides.

\smallskip
If it be possible let the two triangles be constructed so
that
$\left\{\begin{aligned}
  \drawUnitLine{CA}&=\drawUnitLine{DA}\\
  \drawUnitLine{BC}&=\drawUnitLine{BD}
\end{aligned}\right\}$,
then draw \drawUnitLine{CD} and,

\begin{center}
$\drawAngle{C,DCB} = \drawAngle{CDA}$~(pr.~I.5)

$\therefore \drawAngle{DCB} < \drawAngle{CDA}$ and

$\left.\begin{aligned}
  \therefore\drawAngle{DCB} &< \drawAngle{CDA,D}\\
  \text{but~(pr.~I.5)}\enspace\drawAngle{DCB}
    &= \drawAngle{CDA,D}
\end{aligned}\right\}\text{which is absurd,}$
\end{center}

\noindent
therefore the two triangles cannot have their conterminous
sides equal at both extremities of the base.

\begin{flushright}Q.\ E.\ D.\end{flushright}
\end{minipage}

\end{document}
